\chapter{Vorschläge zur Weiterentwicklung} \label{chapter:vorschlaege-weiterentwicklung}

Wie bereits beschrieben, sollte natürlich das Backend der Hauptfokus bei der Weiterentwicklung dieser Webapp sein, damit sie überhaupt als funktionierende Applikation für Ihren angedachten Einsatzzweck verwendet werden kann. Unabhängig vom Backend gibt es nach der zeitlich begrenzten Entwicklung des Frontends auch hier noch einige Ansatzpunkte zur weiteren Verbesserung und für neue Features. Idealerweise sollte außerdem so früh wie möglich auch Feedback von potenziellen Nutzern eingeholt werden, um die Weiterentwicklung dementsprechend zu lenken. Im Folgenden finden sich meine Vorschläge zur Weiterentwicklung der Tracking Editor und Analyse Seiten, alle anderen Teile der Webapp erfüllen meiner Ansicht nach momentan alle notwendigen und wünschenswerten Funktionen.


\subsubsection{Ball Tracking und Ballaktionen}
Die vielversprechendste Ergänzung zu den aktuellen Funktionen wäre meiner Meinung nach die Unterstützung für die Position des Balles und für Ballaktionen wie Pässe, Torschüsse, Freistöße, etc. im Tracking Editor. Dadurch würden sich auch zahllose weitere Möglichkeiten für mehr Analysen ergeben, die im Folgenden näher erläutert werden.

\subsubsection{Unterstützung mehrerer Teams bzw. Gruppen pro Spieler}
Wenn jedem Spieler im Tracking Editor nicht nur sein Team, sondern auch eine beliebige Anzahl an weiteren benutzerdefinierten Gruppen zugeordnet werden könnte, könnte bei der Analyse ebenfalls nach diesen Gruppen unterschieden werden. Mögliche Beispiele im Kontext eines Fußballspiels wären Offensive und Defensive, linke und rechte Seite oder Innenverteidigung und Außenverteidigung. Im Code des Tracking Editors wurde diese Option bereits im Hinterkopf behalten, und die Farbe der jeweiligen Teams bereits so gespeichert, dass in der Zukunft mehrere Sätze an Teams (bzw. dann eher Gruppen) unterstützt werden können. Die Variable \texttt{colorCategoryNumber} in der \texttt{TrackingEditor} Komponente (die momentan nur einen fest zugewiesenen Wert hat) legt fest, welcher Satz an Kategorien verwendet werden soll. Für die Implementierung des beschriebenen Features müsste die \texttt{TrackListItemPlayer} noch so angepasst werden, dass hier jedem Spieler mehrere Gruppen zugeordnet werden können und diese müssten für jeden Spieler in dem ihm zugeordnetem \texttt{bBox} Objekt abgespeichert werden. Die Analyse Seite müsste ebenfalls um eine Auswahl des jeweils gewünschten Satzes an Teams/Gruppen beim Erstellen eines neuen Tiles erweitert werden.

\subsubsection{Spielübergreifende Analysen}

Wenn in der Analyse für jeden Analyse Tile das Spiel ausgewählt werden könnte, könnten verschiedene Spiele miteinander verglichen werden und bestimmte Metriken von einzelnen Spielern, Gruppen von Spielern oder dem gesamten Team über die Saison hinweg verfolgt werden. Dadurch könnten auch eventuell vorhandene Zusammenhänge zwischen der Spielperformance und dem Ausgang des Spiels gefunden werden. Um das zu ermöglichen, müssten allerdings Spieler Ids und Namen über die alle Spiele hinweg konsistent gehalten werden und Metadaten der einzelnen Spiele in Form von teilnehmenden Teams hinzugefügt werden.

\subsubsection{Erweiterte Steuerung im Tracking Editor}

Im Tracking Editor könnte das Bearbeiten der Tracking Daten angenehmer gestaltet bzw. vereinfacht werden, indem ein Zoom implementiert wird. Fabric stellt bereits entsprechende Anknüpfungspunkte zur Verfügung, sodass die Implementierung dieses Features wohl ohne größere Schwierigkeiten möglich sein dürfte. Für eine Anleitung siehe \url{www.fabricjs.com/fabric-intro-part-5#pan_zoom}.\\
Außerdem könnte ein Rechtsklick mit der Maus erkannt und für die Bearbeitung genutzt werden, z.B. indem so verschiedene Bearbeitungstools für die weiter oben genannten Ballevents ausgewählt werden können.\\
Gerade bei hohen Framerates erfordert ein über längere Zeit vorkommender Fehler in beispielsweise der Position eines Spielers sehr viele einzelne Änderungen. Hier wäre es hilfreich, mit Keyframes zu arbeiten, sodass die Position in einem Start- und Endframe manuell korrigiert wird und in allen Zwischenframes linear interpoliert oder vielleicht sogar erneut durch einen Tracking Algorithmus ermittelt wird. Die Wichtigkeit eines solchen Features hängt aber natürlich stark davon ab, wie gut die Tracking Algorithmen bereits funktionieren und ob überhaupt Korrekturen notwendig sind.

\subsubsection{Weitere Analysearten und Freiheiten bei der Gestaltung}

Momentan werden einige grundlegende Arten der Spielanalyse unterstützt. Sobald einige der bereits genannten Weiterentwicklungen erfolgt sind, öffnen diese die Tür für weitere Arten der Spielanalyse, beispielsweise im Fußballkontext Diagramme zum Ballbesitz, zur Anzahl und Verwertung der Chancen oder Standardsituationen, sowie eine Heatmap der Torschusspositionen.\\
Außerdem wäre es denkbar, die Gestaltung der Analyse hinsichtlich Größe und Position der einzelnen Analyse Tiles flexibler zu gestalten, sodass der Nutzer diese mit der Maus verschieben und vergrößern oder verkleinern kann. Gerade bei der Analyse über mehrere Gruppen oder Spiele hinweg wäre es angemessen, die Analyse auch in mehrere Teilbereiche, zum Beispiel in Form von Trennelementen oder verschiedenen umschaltbaren Tabs, gliedern zu können.\\
Zuletzt wäre bei derart umfangreichen Analysen auch die Bereitstellung von Templates sinnvoll. Analog zur aktuellen Bereitstellung von einigen Beispiel Analyse Tiles könnte der Nutzer aus einer Auswahl an Templates schnell eine vorkonfigurierte Analyse erstellen oder sogar eigene Templates abspeichern und für zukünftige Spiele mit wenigen Klicks direkt wiederverwenden. Für die Anpassung der Template Analysen wäre dann eine Funktion zur Bearbeitung der einzelnen Tiles notwendig, die im Moment noch fehlt. Mein Vorschlag wäre, bei den Buttons oben rechts zum Verschieben und Löschen des jeweiligen Tiles noch einen weiteren Button, z.B. mit einem Zahnradsymbol einzubauen. Bei einem Klick würde sich dann die Konfiguration des Tiles erneut öffnen Dafür müsste lediglich der Typ des Tiles zu ``new'' verändert werden (wobei dieser Typ dann sinnvollerweise zu z.B. ``config'' umbenannt werden sollte), damit wieder ein \texttt{TileTypeSelector} mit der aktuellen, bearbeitbaren Konfiguration angezeigt wird. Dafür müsste der Tile, der gerade bearbeitet wird, lediglich seine \texttt{changeTile()} Funktion aufrufen und den neuen Typ sowie alle relevanten, ihm als props übergebenen Argumente als Parameter übergeben. 
    